\chapter{Thua để lớn}
\markboth{Thua để lớn}{Losing to Grow}

\section{Một học sinh rớt đội tuyển rồi mới học cách bền bỉ.}
\begin{truyen}{Một học sinh rớt đội tuyển rồi mới học cách bền bỉ.}{Failure can teach rhythm.}
\chuhoa{E}{m tưởng mình đủ giỏi nên bị loại là cú sốc lớn. Sau lần đó, em không học bốc đồng nữa mà học đều hơn mỗi ngày. Cái thua đầu tiên làm đau, nhưng cũng dạy một nhịp sống mà cái thắng trước đó chưa từng dạy.}
\textit{(The student thought talent was enough, so being rejected was a major shock. After that, study became steady rather than impulsive. The first loss hurt, but it taught a rhythm of life that earlier victories had never taught.)}
\end{truyen}
\begin{baihoc}
Có những thất bại không kéo ta xuống; nó kéo ta vào kỷ luật.
\textit{(Some failures do not pull us down; they pull us into discipline.)}
\end{baihoc}
\begin{ghinhoanh}\textbf{Short English to remember:}

Failure can teach rhythm.
\end{ghinhoanh}
\begin{tuhocnhanh}
1. Cái thua nào đã buộc em sống nghiêm túc hơn? \textit{(Which loss forced you to live more seriously?)}

2. Em đã học được nhịp mới nào từ lần đó? \textit{(What new rhythm did you learn from that event?)}
\end{tuhocnhanh}
\ngancach

\section{Một doanh nhân mất hợp đồng lớn rồi bớt kiêu.}
\begin{truyen}{Một doanh nhân mất hợp đồng lớn rồi bớt kiêu.}{Loss can trim the ego.}
\chuhoa{A}{nh từng nghĩ mình không thể bị thay thế. Khi khách hàng rời đi, anh mới ngồi lại xem cách mình cư xử, cách mình nghe người khác, cách mình coi thường các chi tiết nhỏ. Cái thua làm đau lòng tự ái, nhưng cũng cắt bớt phần kiêu ngạo thừa ra.}
\textit{(He once believed he was irreplaceable. When the client left, he finally reviewed how he behaved, listened, and dismissed small details. The loss hurt his pride, but also cut away some unnecessary arrogance.)}
\end{truyen}
\begin{baihoc}
Nhiều người chỉ bắt đầu dễ chịu hơn sau khi thua một lần thật thấm.
\textit{(Many people only become more bearable after one truly humbling loss.)}
\end{baihoc}
\begin{ghinhoanh}\textbf{Short English to remember:}

Loss can trim the ego.
\end{ghinhoanh}
\begin{tuhocnhanh}
1. Cái tôi của em đã từng bị một thất bại làm mềm lại chưa? \textit{(Has a failure ever softened your ego?)}

2. Sau lần thua đó, em đổi cách cư xử nào? \textit{(What behavior changed after that loss?)}
\end{tuhocnhanh}
\ngancach

\section{Một cô gái chia tay rồi mới học cách đứng bằng hai chân.}
\begin{truyen}{Một cô gái chia tay rồi mới học cách đứng bằng hai chân.}{Loss can return you to yourself.}
\chuhoa{T}{rước đó, cô đặt gần hết niềm vui của mình vào một người. Khi mối quan hệ tan vỡ, cô buộc phải tập sống lại với lịch riêng, bạn bè riêng, sở thích riêng. Cái thua trong tình cảm lại trả cô về với chính cô.}
\textit{(Before, she placed most of her joy in one person. When the relationship ended, she had to rebuild her own schedule, friends, and interests. The loss in love returned her to herself.)}
\end{truyen}
\begin{baihoc}
Có những mất mát không chỉ lấy đi; nó còn buộc ta tìm lại phần mình đã bỏ quên.
\textit{(Some losses do not only take away; they force us to recover the parts of ourselves we neglected.)}
\end{baihoc}
\begin{ghinhoanh}\textbf{Short English to remember:}

Loss can return you to yourself.
\end{ghinhoanh}
\begin{tuhocnhanh}
1. Em có đang đặt quá nhiều bản thân vào một người hay một việc không? \textit{(Are you placing too much of yourself into one person or one thing?)}

2. Nếu mất nó, phần nào của em cần được dựng lại? \textit{(If it were gone, what part of you would need rebuilding?)}
\end{tuhocnhanh}
\ngancach

\section{Một người bị từ chối việc làm rồi mới đi đúng ngành.}
\begin{truyen}{Một người bị từ chối việc làm rồi mới đi đúng ngành.}{A closed door can redirect you.}
\chuhoa{A}{nh buồn vì trượt vị trí mình từng rất muốn. Trong thời gian chững lại, anh thử một hướng khác phù hợp hơn hẳn khả năng và tính cách. Nhiều năm sau, anh biết ơn cú trượt ngày đó vì nếu đậu, có khi anh đã ở sai đường rất lâu.}
\textit{(He was upset after missing a job he badly wanted. During the pause, he tried another path that matched his abilities and personality much better. Years later, he was grateful for that rejection because success there might have kept him on the wrong road.)}
\end{truyen}
\begin{baihoc}
Không phải cánh cửa đóng nào cũng là bất hạnh; có cửa đóng để ta bớt đi sai đường.
\textit{(Not every closed door is a tragedy; some close so we do not stay on the wrong road.)}
\end{baihoc}
\begin{ghinhoanh}\textbf{Short English to remember:}

A closed door can redirect you.
\end{ghinhoanh}
\begin{tuhocnhanh}
1. Có sự từ chối nào về sau em mới thấy là may không? \textit{(Was there any rejection you later saw as fortunate?)}

2. Em có đang tiếc một cánh cửa lẽ ra nên khép không? \textit{(Are you still regretting a door that perhaps needed to close?)}
\end{tuhocnhanh}
\ngancach

\section{Một người thua trong cuộc họp nhưng giữ được sự bình tĩnh.}
\begin{truyen}{Một người thua trong cuộc họp nhưng giữ được sự bình tĩnh.}{Not exploding is also a win.}
\chuhoa{A}{nh bị bác bỏ ý tưởng trước nhiều người. Anh rất khó chịu nhưng không phản ứng nóng. Về sau chính sự điềm tĩnh ấy khiến người khác nể anh hơn cả bản đề xuất ban đầu. Có những lần thua bề mặt lại là thắng lớn về bản lĩnh.}
\textit{(His idea was rejected in front of many people. He felt upset but did not react harshly. Later, that calm response earned more respect than the proposal itself. Some surface losses become deeper victories of character.)}
\end{truyen}
\begin{baihoc}
Không mất mình trong lúc thua đôi khi đã là một cái thắng rất đáng giá.
\textit{(Not losing yourself while losing is sometimes a very valuable victory.)}
\end{baihoc}
\begin{ghinhoanh}\textbf{Short English to remember:}

Not exploding is also a win.
\end{ghinhoanh}
\begin{tuhocnhanh}
1. Khi bị bác bỏ, em thường phản ứng thế nào? \textit{(How do you usually react when being rejected?)}

2. Bản lĩnh nào em muốn giữ trong lần thua tới? \textit{(What kind of composure do you want to keep in your next loss?)}
\end{tuhocnhanh}
\ngancach